%% introduction.tex
%%

%% ==============================
\section{Introduction}
\label{sec:Introduction}
%% ==============================

Infrastructure-as-Code (IaC) misconfigurations -- hardcoded secrets, overly permissive file modes, unencrypted protocols -- are a recurring cause of real security incidents. Rule-based scanners like Checkov catch well-known, enumerable patterns reliably, and recent work has shown that semantics-aware ML models (fine-tuned on both code and accompanying natural-language comments) can substantially outperform older statistical baselines at this task.

War et al. \cite{War25} demonstrate exactly this with CodeBERT and LongFormer, but their own ablation study exposes a specific weakness: when natural-language context is stripped away, precision collapses while recall holds or rises, meaning the model starts flagging large amounts of safe code as misconfigured. This matters in practice, since real-world IaC repositories are often sparsely commented, unlike the curated, comment-rich datasets used to train and evaluate their models.

\textbf{Research question:} can a hybrid detector -- combining a comment-independent rule layer with the same ML classifier, gated more conservatively -- recover the precision lost when natural-language context is unavailable, and if so, at what cost?

\textbf{Objectives:}
\begin{itemize}
  \item Reproduce the paper's ablation finding (precision collapse without natural-language context) on a controlled synthetic benchmark.
  \item Design the benchmark so that some misconfigurations are code-alone-solvable ("easy") and some are only resolvable via a natural-language comment ("hard"), making the reproduction genuine rather than trivially avoidable.
  \item Build a hybrid rule+ML detector targeting the gap, and report its actual trade-offs rather than only the metrics that look favourable.
\end{itemize}

The remainder of this report is structured as follows: Section~\ref{sec:rw} reviews the baseline paper and the gap it names; Section~\ref{sec:Design} describes the easy/hard benchmark design and the three detectors compared; Section~\ref{sec:Implementation} describes the codebase and deployment; Section~\ref{sec:Evaluation} reports results across 5 seeds and a threshold sweep; Section~\ref{sec:Conclusion} discusses the trade-off and future work.
